\section{The two cones and the void}

Owner ruling 2026-09-21, ``the admission boundary, restored.'' A charged
valuation is bracketed at every cell between two fixed points of the settling
closure: the unit cone, the least the given evidence forces, and the counit
cone, the most it allows. The gap between them, cell by cell, is the void ---
under-determination that is real because the carrier is measured
(\cref{thm:charged-measured}), not an artifact of where a value happened to
be stamped. \texttt{CValuation} abbreviates \(\texttt{PredicationCell} \to
\texttt{Charged}\); the cones below act pointwise, cell by cell, through
\texttt{Charge.settle} on each cell's own value.

\begin{definition}[The unit cone: settle]\label{thm:cvaluation-settle}\lean{Obsmetrologia.CValuation.settle}\leanok\uses{thm:closure-operator, thm:settling, thm:predication}
\(\texttt{settle}\,c\,v\) is the valuation obtained by closing each cell's
charged value under the closure \(c\), pointwise:
\(\varphi \mapsto \texttt{Charge.settle}\,c\,(v\,\varphi).\mathrm{value}\).
This is the least the evidence in \(v\) forces: no free side is raised before
closing.
\end{definition}

\begin{definition}[The counit cone: brim]\label{thm:cvaluation-brim}\lean{Obsmetrologia.CValuation.brim}\leanok\uses{thm:cvaluation-settle, thm:charged-presumefree-ge}
\(\texttt{brim}\,c\,v\) first presumes every evidence-free side of each
cell's charged value to present (\cref{thm:charged-presumefree-ge}), then
closes under \(c\), pointwise. This is the most the evidence in \(v\)
allows: every side not pinned by a witness is raised before closing.
\end{definition}

\begin{theorem}[settle \(\le_k\) brim]\label{thm:cvaluation-settle-le-brim}\lean{Obsmetrologia.CValuation.settle_le_brim}\leanok\uses{thm:cvaluation-settle, thm:cvaluation-brim}
At every cell, \(\texttt{settle}\,c\,v\,\varphi \le_k \texttt{brim}\,c\,v\,
\varphi\): the unit cone never sits above the counit cone. This follows from
\(c\)'s monotonicity together with \cref{thm:charged-presumefree-ge}'s
inflationary presumption --- the least the evidence forces can never exceed
the most it allows.
\end{theorem}

\begin{definition}[The void]\label{thm:cvaluation-void}\lean{Obsmetrologia.CValuation.Void}\leanok\uses{thm:cvaluation-settle-le-brim}
A cell \(\varphi\) is \emph{void} under \(c\) and \(v\) when the two cones
disagree there: \(\texttt{settle}\,c\,v\,\varphi \neq \texttt{brim}\,c\,v\,
\varphi\). The void is the set of cells the closure pins from neither side
--- the monograph's reading ``the void is space.''
\end{definition}

\begin{proposition}[Void, restated as a strict gap]\label{thm:cvaluation-void-iff}\lean{Obsmetrologia.CValuation.void_iff}\leanok\uses{thm:cvaluation-void}
A cell is void exactly when the counit cone is \textbf{not} \(\le_k\) the
unit cone there: since \cref{thm:cvaluation-settle-le-brim} already gives one
direction always, disagreement is equivalent to the reverse inequality
failing.
\end{proposition}

\begin{theorem}[Fully evidenced cells are never void]\label{thm:cvaluation-not-void-of-full}\lean{Obsmetrologia.CValuation.not_void_of_full}\leanok\uses{thm:cvaluation-void-iff, thm:charged-presumefree-ge}
If a cell's evidence is non-empty on both sides, it is not void under any
closure: with nothing left free to presume, \texttt{brim} and \texttt{settle}
agree there pointwise (\cref{thm:charged-presumefree-ge}'s companion fact,
that presumption changes nothing once every side is already witnessed).
\end{theorem}

\begin{theorem}[A bare articulated cell is void]\label{thm:cvaluation-articulate-void}\lean{Obsmetrologia.CValuation.articulate_void}\leanok\uses{thm:cvaluation-void, thm:cell-articulate-corner, thm:charged-presumefree-articulate}
Reading the constant valuation of a freshly articulated cell under the
identity-seeded join closure, the claim side is void: \texttt{settle} leaves
the cell at potential while \texttt{brim} presumes it to stable
(\cref{thm:charged-presumefree-articulate}), and the two disagree. This is
the pointwise witness of the slack \cref{thm:knowledge-least} always carries
at an un-engaged claim.
\end{theorem}

\begin{theorem}[A cancelled conflict is not void]\label{thm:cvaluation-cancelled-not-void}\lean{Obsmetrologia.CValuation.cancelled_not_void}\leanok\uses{thm:cvaluation-void, thm:charged-meetk-not-measured}
Reading the constant valuation of the knowledge meet of an articulated cell
and an active cell built by fiat, the resulting cell is \textbf{not} void:
both sides are witnessed (\cref{thm:charged-meetk-not-measured}'s meet, value
driven to \(\bot_k\) with both witnesses kept), so there is nothing left free
for \texttt{brim} to raise, and the two cones agree. This is the formal
separation the theory needs between ``a cell no evidence has reached'' (void)
and ``a cell whose evidence has reconciled to a cancellation'' (not void,
even though both currently read \(\bot_k\)).
\end{theorem}
